\chapter{Matcheo de patrones temporales}
\label{ch:matcher}

El matcher es el hot path del motor. Cada hunt que corre el
analista, cada llamada \code{run\_hunt} desde Claude, cada
invocación CLI termina acá. La implementación trae tres variantes
coordinadas:

\begin{itemize}
    \item \textsc{Search-Temporal-Pattern} --- el entry point;
        despacha trabajo a workers rayon y junta resultados.
    \item \textsc{Dfs-Match-Iterative} --- el DFS iterativo base
        que se ejecuta por vértice de partida.
    \item \textsc{Dfs-Match-Smart} --- variante con cota de
        anomalía que mantiene un min-heap top-$K$ y poda ramas
        cuyo estimado corriente las hace inadmisibles.
\end{itemize}

Un cuarto camino, \textsc{Search-Temporal-Pattern-Simd}, delega
en \code{libgraphmatch} o en el módulo puro-Rust \code{simd\_rust}
(Capítulo~\ref{ch:simd}) y está gateado por el feature \code{simd}
más una prueba de CPU en runtime.

\section{Estado compartido}

\begin{lstlisting}[style=rust]
// Scratch alocado una vez por worker y reusado entre starts.
struct DfsFrame {
    sid:          StrId,
    step_idx:     usize,
    edge_hi:      usize,   // cota superior de aristas salientes
    edge_cursor:  usize,   // siguiente indice de arista a probar
}

struct SmartDfsFrame {
    sid:              StrId,
    step_idx:         usize,
    edge_hi:          usize,
    edge_cursor:      usize,
    path_anomaly_sum: f64,   // suma de node_anomaly_estimate a lo largo del path
}
\end{lstlisting}

Los frames de la pila son structs comunes guardados en un
\code{Vec<Frame>} que se reusa entre vértices de partida. Este es
el ahorro de alocación más grande del matcher --- la
implementación recursiva anterior alocaba por frame en la pila
del sistema y no podía reiniciarse trivialmente.

\begin{invariant}[$k$-simplicidad]
El \code{visit\_count: HashMap<StrId, usize>} corriente mantenido
dentro de cada caminata DFS satisface \code{visit\_count[v] <= k}
para todo $v$ en el camino actual, donde $k$ es
\code{hypothesis.k\_simplicity}. Violarlo admitiría caminos que
revisitan un vértice más de $k$ veces, contradiciendo la
semántica de la hipótesis. El algoritmo impone el invariante en
el push: si incrementar \code{visit\_count[dest]} excedería $k$,
la arista se descarta.
\end{invariant}

\section{Entry point: \textsc{Search-Temporal-Pattern}}

El algoritmo externo particiona el trabajo por vértice de
partida. Por cada entidad cuyo tipo matchea el primer paso de
$H$, un thread rayon ejecuta \textsc{Dfs-Match-Iterative} sobre
ese vértice con su propia pila scratch y su propio buffer de
resultados. Cuando todos los workers terminan, los resultados se
mergean, puntúan, deduplican y paginan
(Capítulo~\ref{ch:dedup-pagination}).

\begin{algorithm}
\caption{\textsc{Search-Temporal-Pattern}}\label{alg:search-entry}
\KwIn{grafo $G$, hipótesis $H$, ventana opcional $(t_0, t_1)$,
      cap opcional $\mathrm{max}$}
\KwOut{vector de vectores de caminos}
$\mathrm{starts} \gets $ \code{G.type\_index[H.steps[0].origin\_type]}\;
\If{$\mathrm{starts} = \emptyset$}{\Return $\emptyset$\;}
$\mathrm{cap} \gets $ \code{AtomicUsize(0)}\;
$\mathrm{results} \gets $ map paralelo sobre $\mathrm{starts}$:
    \ForEach{$s \in \mathrm{starts}$ en paralelo}{
        $\mathrm{local} \gets $ \textsc{Dfs-Match-Iterative}($G, H, s, t_0, t_1, \mathrm{cap}, \mathrm{max}$)\;
        \Return $\mathrm{local}$\;
    }
\Return concat de $\mathrm{results}$\;
\end{algorithm}

Cada worker lee el \code{AtomicUsize} cap en cada push de
candidato, de modo que cuando el conteo agregado alcanza
$\mathrm{max}$, los workers restantes paran lo antes posible. Es
un patrón clásico de early-exit: el count puede rebasar
$\mathrm{max}$ momentáneamente porque los workers \emph{commit}
y luego chequean, pero el overshoot está acotado por la cantidad
de workers rayon (típicamente 8--16).

\section{El DFS iterativo}

\begin{algorithm}
\caption{\textsc{Dfs-Match-Iterative}}\label{alg:dfs-iter}
\KwIn{grafo $G$, hipótesis $H$, start $s_0$, ventana $(t_0, t_1)$,
      contador compartido $\mathrm{cap}$, cota $\mathrm{max}$}
\KwOut{lista de caminos matcheantes rooteados en $s_0$}
$\mathrm{results} \gets \emptyset$;
$\mathrm{stack} \gets \emptyset$;
$\mathrm{path} \gets [s_0]$;
$\mathrm{visit} \gets \{s_0 : 1\}$\;
\textsc{Push-Edges-For-Step}($\mathrm{stack}$, $G$, $s_0$, paso $0$, $t_0, t_1$)\;
\While{$\mathrm{stack} \neq \emptyset$}{
    $f \gets $ tope de $\mathrm{stack}$\;
    \If{$f.\text{edge\_cursor} \geq f.\text{edge\_hi}$}{
        popear $f$;
        popear último vértice de $\mathrm{path}$;
        decrementar $\mathrm{visit}[\cdot]$;
        \Continue\;
    }
    $c \gets $ \code{G.edge\_store[f.sid][f.edge\_cursor]}\;
    $f.\text{edge\_cursor} \gets f.\text{edge\_cursor} + 1$\;
    \tcp{filtro: tipo de relación}
    \If{$c.\text{rel\_type\_tag} \neq $ paso[$f.\text{step\_idx}$].rel}{
        \Continue\;
    }
    \tcp{filtro: tipo de entidad destino vía tag SoA}
    $u \gets c.\text{dest\_sid}$\;
    \If{$\code{G.entity\_type\_tags}[u.0] \neq $ paso[$f.\text{step\_idx}$].dst}{
        \Continue\;
    }
    \tcp{filtro: ventana temporal}
    \If{$c.\text{timestamp} \notin [t_0, t_1]$}{\Continue\;}
    \tcp{filtro: predicados de arista (sobre metadata)}
    \If{\textsc{Evaluate-Predicates}(paso.edge\_predicates, $c$, $G$) = falso}{
        \Continue\;
    }
    \tcp{prevención de ciclos vía $k$-simplicidad}
    \If{$\mathrm{visit}[u] + 1 > H.k$}{\Continue\;}
    $\mathrm{visit}[u] \gets \mathrm{visit}[u] + 1$;
    anexar $u$ a $\mathrm{path}$\;
    \If{$f.\text{step\_idx} + 1 = $ len($H.\text{steps}$)}{
        registrar $\mathrm{path}$ en $\mathrm{results}$\;
        incrementar $\mathrm{cap}$ atómicamente;
        \If{$\mathrm{cap} \geq \mathrm{max}$}{salir del loop\;}
        $\mathrm{visit}[u] \gets \mathrm{visit}[u] - 1$;
        popear último vértice;
        \Continue\;
    }
    empujar frame nuevo: $(u, f.\text{step\_idx}+1, \ldots)$ sobre $\mathrm{stack}$\;
    \textsc{Push-Edges-For-Step}($\mathrm{stack}$, $G$, $u$, $f.\text{step\_idx}+1$, $t_0, t_1$)\;
}
\Return $\mathrm{results}$\;
\end{algorithm}

\textsc{Push-Edges-For-Step} es una rutina de una línea que hace
binary search en \code{edge\_store[v]} para la primera arista con
\code{timestamp} $\geq t_0$ y setea \code{edge\_hi} a la primera
arista con \code{timestamp} $> t_1$. El sort por timestamp que se
hace en finalize (Capítulo~\ref{ch:graph-hunter}) es lo que
justifica que la binary search funcione.

\begin{theorem}[Corrección del matcher]
Al retornar, \textsc{Dfs-Match-Iterative} produce exactamente el
conjunto de caminos temporales de longitud $\ell$ rooteados en
$s_0$ que satisfacen la hipótesis $H$ dentro de $[t_0, t_1]$ y
respetan la $k$-simplicidad.
\end{theorem}

\begin{proof}[Demostración (esbozo)]
Por inducción sobre la profundidad de la hipótesis. El caso base
($\ell = 1$) es el primer push en \code{stack}; el chequeo es
paso a paso equivalente a las restricciones de la hipótesis sobre
la primera arista. El paso inductivo observa que cada arista
admitida extiende el camino exactamente un vértice, actualiza
\code{visit} consistentemente, y que el chequeo de ciclos rechaza
exactamente aquellas extensiones que excederían $k$-simplicidad.
La pila se desenrolla restaurando \code{visit} y el camino, de
modo que cada rama se explora independientemente; ningún camino
se omite por backtracking faltante. \qed
\end{proof}

\begin{complexity}
El peor caso es un grafo donde toda arista matchea todo paso.
Entonces desde un start $s_0$ el DFS visita hasta
$\bar{b}^\ell$ caminos (con $\bar{b}$ = grado saliente promedio
de aristas calificantes por paso y $\ell$ = cantidad de pasos),
dando $\Ord{\bar{b}^\ell \cdot \ell}$ por start. Cruzando los
$n_s$ candidatos de partida, el total es
\[
    T_{\text{match}}(G, H) = \Ord{n_s \cdot \bar{b}^\ell \cdot \ell}.
\]
El espacio por worker es $\Ord{\ell}$ para la pila y el camino más
$\Ord{\ell \cdot k}$ para el mapa \code{visit} (cada entrada es un
vértice, y el camino tiene a lo sumo $\ell$ vértices distintos cada
uno visitado a lo sumo $k$ veces). Con $W$ workers, el espacio
total del matcher es $\Ord{W \cdot \ell \cdot k}$. En la práctica
$\bar{b}$ es diminuto porque el fast-path por tag de tipo elimina
la mayoría de las aristas salientes en la primera comparación de
byte; mediciones empíricas desde
\code{benches/hunt\_latency.rs} están en el Ap.~\ref{app:bench}.
\end{complexity}

\section{La variante smart}

\textsc{Dfs-Match-Smart} se usa cuando el llamador pide los top-$K$
caminos más anómalos en vez de ``todos los caminos''. Difiere del
iterativo sólo en el orden de candidatos y el acumulador de
resultados:

\begin{itemize}
    \item \textbf{Acumulador}: un min-heap de tamaño a lo sumo $K$
        ordenado por puntaje compuesto de anomalía. Cuando está
        lleno, los caminos nuevos se aceptan si y sólo si puntúan
        más alto que el mínimo del heap.
    \item \textbf{Cota corriente}: cada frame trackea
        \code{path\_anomaly\_sum}, la suma de
        \code{node\_anomaly\_estimate} (Capítulo~\ref{ch:anomaly})
        a lo largo del camino hasta el momento. Extender el
        camino no puede decrecer esta suma; si la suma ya cae por
        debajo del mínimo del heap, la rama se poda.
\end{itemize}

\begin{algorithm}
\caption{\textsc{Merge-Into-TopK}}\label{alg:merge-topk}
\KwIn{heap local $L$, heap global $G$, capacidad $K$}
\ForEach{$p \in L$}{
    \If{$|G| < K$}{$G$.push($p$)}
    \ElseIf{$\text{score}(p) > \min(G)$}{
        $G$.pop\_min(); $G$.push($p$);
    }
}
\end{algorithm}

\begin{complexity}
$\Ord{(|L| + |G|) \log K}$. Cuando el matcher junta los workers
rayon, cada heap local de tamaño $\leq K$ se mergea en un único
heap global a costo total $\Ord{W K \log K}$, despreciable frente
al trabajo del DFS.
\end{complexity}

\section{Paralelismo}

El matcher es embarazosamente paralelo en el vértice de partida.
Como cada DFS de worker lee el grafo exclusivamente (todos los
escritores están bloqueados en un \code{RwLock::write} a nivel
sesión), no hay sincronización entre workers más allá del
contador compartido \code{cap} y la reducción final en un único
vector de resultados.

\begin{theorem}[Speedup paralelo]
Dados $W$ workers rayon y $n_s$ vértices de partida con carga de
DFS aproximadamente igual por start, el matcher alcanza un
speedup de $\min(W, n_s)$ hasta el costo constante de la
reducción final.
\end{theorem}

La prueba es el teorema de Brent aplicado a la partición trivial:
el DFS de cada start no tiene dependencias entre starts, y la
reducción final es $\Ord{\text{total resultados}}$, típicamente
órdenes de magnitud menor que el DFS en sí.

\begin{inthecode}
\filepath{graph\_hunter\_core/src/graph.rs:461} ---
\code{search\_temporal\_pattern}.\\
\filepath{graph\_hunter\_core/src/graph.rs:1121} ---
\code{dfs\_match\_iterative}.\\
\filepath{graph\_hunter\_core/src/graph.rs:782} ---
\code{search\_temporal\_pattern\_smart}.\\
\filepath{graph\_hunter\_core/src/graph.rs:932} ---
\code{dfs\_match\_smart}.\\
\filepath{graph\_hunter\_core/benches/hunt\_latency.rs} ---
benchmarks de performance.
\end{inthecode}
